\section{Related Work}
\label{sec:related}

Orbital inspection draws on three bodies of literature: relative spacecraft
motion, robotic view planning, and middleware-supported simulation. The
Clohessy--Wiltshire equations remain a common starting point for rendezvous
and proximity operations near a circular chief orbit because they provide a
linear, interpretable model of relative translation \cite{clohessy1960terminal}.
Although high-fidelity perturbation models are required for flight design, a
CW/HCW baseline is appropriate for early algorithmic comparisons because it
makes assumptions explicit and permits repeatable propagation, control, and
transfer-cost calculations \cite{wertz2011space}.

Robotic inspection and next-best-view planning study how a robot should select
sensor poses that reconstruct or inspect an object efficiently. Classical
surveys emphasize visibility, occlusion, sensor range, view quality, and the
trade-off between coverage and motion cost \cite{scott2003view}. Orbital
inspection adds constraints that are less prominent in terrestrial settings:
viewpoints must be dynamically reachable under relative orbital motion, safety
margins must be maintained around delicate structure, and fuel consumption is a
primary mission metric. This motivates planners that do not choose views solely
by information gain, but also by transfer feasibility and clearance.

Simulation frameworks for space robotics must also be treated carefully. A
rendered Gazebo scene can show useful visual context, but standard rigid-body
physics in a robotics simulator is not a substitute for an orbital relative
dynamics model. OrbInspect therefore uses Gazebo Harmonic for visualization of
the ISS model and chaser while the ROS~2 dynamics node publishes the chaser
state used by planning, control, safety, and logging. This separation follows
the philosophy that numerical claims should come from logged planning and
execution data, while visualization should support interpretation and
interface validation.

The closest methodological position of this work is therefore not a single
new optimizer in isolation, but an integrated, reproducible evaluation stack:
coverage-aware viewpoint planning, CW-feasible trajectory generation, safety
monitoring, ROS~2 execution, and Gazebo visualization are connected through
recorded artifacts that can be audited after each run.
